\section{Background and Model}
\label{sec:background}

This section reviews level-synchronous and direction-optimizing BFS, defines the progress guarantees we use, and states the model under which our method is wait-free. Throughout the paper, we say \emph{vertex} for an element of a graph and \emph{node} for an element of a linked list.

\subsection{Breadth-First Search}
\label{sec:bg-bfs}

We treat every graph as undirected. A graph $G = (V, E)$ has $n$ vertices, numbered $0$ to $n-1$, and we store each of its edges $\{u, v\}$ as the two directed edges $(u, v)$ and $(v, u)$; $m$ counts these directed edges. The graph is kept in compressed sparse row (CSR) form: the neighbors of a vertex $u$ are $\mathit{adj}[\mathit{offsets}[u]]$ up to, not including, $\mathit{adj}[\mathit{offsets}[u+1]]$, and the \emph{degree} of $u$ is $\mathit{offsets}[u+1] - \mathit{offsets}[u]$.

A BFS from a source vertex $s$ computes the distance $\delta(v)$ of every vertex $v$: the number of edges on a shortest path from $s$ to $v$, or $-1$ if there is no such path. Level $d$ is the set $L_d = \{v : \delta(v) = d\}$, so $L_0 = \{s\}$. The search keeps the distances it has found in an array $\mathit{dist}$, which holds $-1$ for every vertex at the start. A vertex is \emph{discovered} when its distance is set, and \emph{expanded} when its neighbors are examined; the level being expanded is the \emph{frontier}.

A top-down step expands the frontier $L_d$: it examines every neighbor $v$ of every vertex in $L_d$, and discovers each $v$ that is still undiscovered, setting $\mathit{dist}[v] = d+1$ and adding $v$ to $L_{d+1}$. A level-synchronous parallel BFS~\cite{bader_bfs, charles_bfs} divides each frontier among the threads, as \textsc{TopDownStep} in Algorithm~\ref{alg:do-bfs} shows. When several threads find the same undiscovered vertex, a compare-and-swap (CAS) on its distance lets only one of them add it to the next level. The step ends at a barrier, since no thread may expand $L_{d+1}$ before all of it has been discovered: otherwise a thread could reach a vertex of $L_{d+1}$ that is still undiscovered from another vertex of $L_{d+1}$, and give it the distance $d+2$.

A bottom-up step~\cite{beamer_dobfs} works the other way around: every undiscovered vertex $v$ scans its own neighbors for one in the frontier, and if it finds one, it sets $\mathit{dist}[v] = d+1$, adds $v$ to $L_{d+1}$, and stops scanning (\textsc{BottomUpStep} in Algorithm~\ref{alg:do-bfs}). Only the thread that scans $v$ writes its distance, so the step needs no CAS, but it needs a fast test of whether a vertex is in the frontier; GAPBS keeps the frontier as a bitmap for this test. When the frontier is large, a bottom-up step examines far fewer edges than a top-down step: most undiscovered vertices find a frontier neighbor among their first few neighbors, while a top-down step examines every edge of the frontier, most of which lead to vertices that are already discovered.

A direction-optimizing BFS chooses the direction of every step. The rule of Beamer et al., as GAPBS implements it, keeps two counts: $m_f$, the number of edges of the frontier, which is the sum of the degrees of its vertices, and $m_u$, an estimate of the edges not yet examined, which starts at $m$ and drops by $m_f$ with every top-down step. A search going top-down switches to bottom-up once $m_f > m_u / \alpha$. A search going bottom-up continues while the frontier grows or holds more than $n / \beta$ vertices, and switches back to top-down otherwise. GAPBS sets $\alpha = 15$ and $\beta = 18$; our method uses the same rule and values (Section~\ref{sec:dir-opt}).

\begin{algorithm}
\caption{Level-synchronous direction-optimizing BFS, as in GAPBS. Every step ends at a barrier.}
\label{alg:do-bfs}
\begin{algorithmic}[1]
\Procedure{TopDownStep}{$F$, $d$}
    \State $F' \gets \emptyset$
    \ForAll{$u \in F$ \textbf{in parallel}}
        \ForAll{neighbors $v$ of $u$}
            \If{$\mathit{dist}[v] = -1$ \textbf{and} \Call{CAS}{$\mathit{dist}[v], -1, d+1$}}
                \State add $v$ to $F'$
            \EndIf
        \EndFor
    \EndFor
    \State \textbf{barrier} \Comment{waits for the slowest thread}
    \State \Return $F'$
\EndProcedure
\Statex
\Procedure{BottomUpStep}{$F$, $d$}
    \State $F' \gets \emptyset$
    \ForAll{$v \in V$ with $\mathit{dist}[v] = -1$ \textbf{in parallel}}
        \ForAll{neighbors $u$ of $v$}
            \If{$u \in F$}
                \State $\mathit{dist}[v] \gets d+1$; add $v$ to $F'$
                \State \textbf{break}
            \EndIf
        \EndFor
    \EndFor
    \State \textbf{barrier} \Comment{waits for the slowest thread}
    \State \Return $F'$
\EndProcedure
\Statex
\Procedure{BFS}{$s$}
    \State $\mathit{dist}[s] \gets 0$; $F \gets \{s\}$; $d \gets 0$
    \While{$F \neq \emptyset$}
        \If{the direction rule chooses top-down}
            \State $F \gets$ \Call{TopDownStep}{$F$, $d$}
        \Else
            \State $F \gets$ \Call{BottomUpStep}{$F$, $d$}
        \EndIf
        \State $d \gets d+1$
    \EndWhile
\EndProcedure
\end{algorithmic}
\end{algorithm}

\subsection{Progress Guarantees}
\label{sec:bg-progress}

Progress guarantees are properties of methods~\cite{aomp}. A method is \emph{blocking} if a delay of one thread can keep the calls of other threads from finishing, as when a call waits for a lock that a delayed thread holds, or at a barrier that a delayed thread has not reached. A method is \emph{lock-free} if it guarantees that infinitely often some call finishes in a finite number of steps: the system as a whole makes progress, although a particular call may never finish. A method is \emph{wait-free} if every call finishes in a finite number of the calling thread's own steps, whatever the other threads do~\cite{herlihy_waitfree}, and \emph{bounded wait-free} if there is a bound on that number. A call of a wait-free method therefore never waits for another thread, and even a thread that stops for good delays no other. The usual way to make a method wait-free is \emph{helping}: instead of waiting for a delayed thread to finish its part of the work, a thread finishes that part itself~\cite{aomp, kogan_fastwf}.

The BFS method of Algorithm~\ref{alg:do-bfs} is blocking: a thread that is delayed before the barrier of a step delays every other thread there, and one that stops for good keeps them waiting forever.

\subsection{System Model}
\label{sec:bg-model}

There are $T$ threads, numbered $0$ to $T-1$, which communicate through shared memory. The number of threads is fixed before a search, and every thread knows it and its own number. Threads are asynchronous: any thread can be delayed for any length of time between any two of its steps, or stop for good, and we assume nothing about their relative speeds. The graph does not change during a search.

Memory accesses follow total store order (TSO), the memory model of x86 processors~\cite{x86tso, intelmanual, amdmanual}. Aligned loads and stores of up to eight bytes are atomic, and so is a compare-and-swap (CAS) on such a location. The stores of a thread become visible to all other threads at the same time, and in the order in which the thread made them, and the loads of a thread take effect in program order. The one reordering that TSO allows is that a load can take effect before an earlier store of the same thread to another location becomes visible, while the store still waits in the thread's store buffer. Our method relies only on the order between two stores of a thread and between two loads of a thread, and on CAS, which x86 performs as a locked instruction that also orders the accesses around it. We write the method with ordinary loads and stores, and put compiler barriers where their order matters, so that the compiler keeps those accesses in program order. On a weaker memory model, the stores that must stay ordered would need release semantics, and the loads that read them acquire semantics.

Each thread allocates the memory it needs during a search from its own arena, which hands out blocks by moving a pointer and asks the system for more memory only when it runs out. We assume that a thread can allocate memory in a bounded number of its own steps.

\subsection{The BFS Method}
\label{sec:bg-method}

Every thread calls the method \textsc{BFS}$(s)$ once per search, with the same source $s$. Before the first call, $\mathit{dist}[v] = -1$ for every vertex $v$, and the method's other per-vertex arrays hold their initial values. Resetting these arrays between searches is not part of the method; since a search reuses them, the next search starts only once every call of the previous one has returned.

When any call returns, $\mathit{dist}[v] = \delta(v)$ for every vertex $v$, and no step that any thread takes afterwards changes $\mathit{dist}$. A search is therefore complete when its first call returns, and that is when we stop timing it (Section~\ref{sec:experiments}). Section~\ref{sec:proof} proves these properties, and that every call returns after a bounded number of the calling thread's own steps.
